\section{TinyRL --- Community OSS Framework}

\smallskip\noindent\textit{Study Guide companion: Study Guide Section 8 (Research Proposals --- Proposal 2) covers the interview-ready version of this material.}
\smallskip

\begin{whythis}
RLVR (Reinforcement Learning with Verifiable Rewards) has replaced RLHF as the dominant paradigm for frontier model training. Models are trained for weeks against deterministic verifiers --- unit tests, math checkers, formal proof systems --- and the failure modes that emerge (reward hacking, entropy collapse, verifier overfitting) are poorly understood even by practitioners. Yet every paper studying these phenomena requires standing up a custom training stack, designing bespoke environments, and running experiments that take days on expensive hardware.

TinyRL fills the same role MNIST filled for computer vision: a standardized, minimal starting point that everyone can benchmark against. MNIST did not replace ImageNet. It gave researchers a shared substrate where ideas could be tested cheaply, compared cleanly, and built upon publicly. TinyRL does this for RLVR. The headline environment, TinyRL-Code, is RLVR in miniature --- verifiable rewards, deterministic unit tests, emergent strategies, and a reward hacking gallery --- all on a single GPU in under an hour.

As a community asset, TinyRL creates network effects. Contributors add environments, researchers compare results across algorithms, the reward hacking gallery attracts attention from practitioners who have seen the same exploits at scale. The barrier to entry for RLVR experimentation drops from ``access to a frontier training cluster'' to ``a laptop and 30 minutes.'' That is the lever.
\end{whythis}

\subsection{Design Principles}

The TinyStories paper (arXiv 2305.07759) showed that language model phenomena --- memorization, generalization, induction heads, grokking --- are observable at toy scale if the environment is designed carefully. Four properties made it work:

\begin{enumerate}
  \item \textbf{Controlled data distribution.} Stories were generated from a fixed vocabulary and grammar, so the input distribution was known and manipulable.
  \item \textbf{Tasks learnable by small models.} GPT-2 scale (124M parameters) was sufficient to exhibit the phenomena of interest.
  \item \textbf{Clean evaluation.} Human raters could assess story quality reliably; there was no ambiguity about what ``good'' meant.
  \item \textbf{Phenomena of interest observable at small scale.} Memorization curves, phase transitions, emergent capabilities --- all visible without scaling to billions of parameters.
\end{enumerate}

TinyRL maps each of these to RLVR:

\begin{enumerate}
  \item \textbf{Curated prompt sets.} A fixed bank of 50--100 programming puzzles (or IMDB prompts, or quantization tasks), sampled reproducibly. The training distribution is known.
  \item \textbf{GPT-2 scale models.} 124M--350M parameters. Cheap to train, cheap to share, easy to ablate.
  \item \textbf{Deterministic verification.} Unit tests pass or fail; sentiment scores are computed by a fixed classifier; inference cost is measurable. No human raters, no ambiguity.
  \item \textbf{RLVR phenomena visible in under one hour.} Reward hacking, entropy collapse, emergent chain-of-thought under verifiable pressure, KL divergence growth --- all observable on a single GPU before lunch.
\end{enumerate}

\begin{thinkfirst}{What Makes a Good Minimal RL Environment?}
Before reading further, design the simplest RLVR environment you can imagine that still exhibits \textit{verification hacking}. Specify: the model size, the task format, the verifier, and the exploit you predict the model will find. How would you know the model is gaming the verifier rather than solving the task? Write this down before continuing.
\end{thinkfirst}

\subsection{RLVR-First Framing}

RLVR trains a policy against rewards computed by an objective, external verifier --- a function that takes a completion and returns a scalar without human involvement. Math correctness is verified by a symbolic checker. Code correctness is verified by unit tests. The reward is not a neural proxy; it is a ground-truth signal about whether the model solved the problem.

This matters because RLHF's failure mode is reward model gaming: the policy learns to produce text that scores highly on the reward model without actually being better. RLVR's promise is that the verifier is non-gameable in principle --- if the unit tests pass, the code is correct. In practice, the verifier is gameable, and the exploits are more subtle: hardcoded outputs, format manipulation, specification gaps. TinyRL-Code is designed to make these exploits visible.

TinyRL has three environments. \textbf{TinyRL-Code} is the headline RLVR environment: unit-test-verified programming puzzles. \textbf{TinyRL-Align} is the RLHF baseline: sentiment optimization against a neural reward model. \textbf{TinyRL-Quant} is a RLVR-adjacent environment: reward is fully computable (accuracy $\times$ efficiency) but the action space is discrete and combinatorial. The contrast between Code and Align is itself the core pedagogical payload --- it shows concretely why the field moved from neural proxy rewards to verifiable rewards, and at what cost.

\subsection{Framework Architecture}

\subsubsection{Standardized Training Loop}

TinyRL exposes a single interface that every algorithm implements against:

\begin{itemize}
  \item \texttt{env.reset()} --- returns a prompt string sampled from the environment's prompt bank.
  \item \texttt{env.step(completion)} --- takes a string completion, runs the verifier, returns a scalar reward.
  \item \texttt{env.eval(policy)} --- runs the policy on the held-out evaluation set and returns environment-specific metrics.
\end{itemize}

Algorithms (PPO, GRPO, REINFORCE, DPO) are implemented as wrappers around this interface. Swapping algorithms requires changing one line: \texttt{trainer = PPOTrainer(env)} becomes \texttt{trainer = GRPOTrainer(env, G=8)}. The same environment definition works with every algorithm, enabling clean cross-algorithm comparisons.

\subsubsection{Metrics Dashboard}

Every training run automatically logs:

\begin{itemize}
  \item \textbf{Reward curve} --- mean reward per rollout batch, smoothed and raw.
  \item \textbf{KL divergence} --- $D_{\mathrm{KL}}(\pi_\theta \| \pi_{\mathrm{ref}})$ computed on each batch; the primary indicator of policy drift.
  \item \textbf{Policy entropy} --- $H(\pi_\theta)$; collapse signals mode collapse and reward overfitting.
  \item \textbf{Generation quality samples} --- 5 random completions logged at each checkpoint for qualitative inspection.
  \item \textbf{Gradient norm} --- $\ell_2$ norm of the gradient before clipping; instability here precedes training collapse.
\end{itemize}

All metrics are logged in a common schema (Weights \& Biases or TensorBoard), making cross-environment and cross-algorithm comparisons trivial. The leaderboard ingests the same metrics.

\subsubsection{Pluggable Environments}

Each environment is a single Python file defining three objects:

\begin{itemize}
  \item \texttt{prompts} --- a list of strings (or a callable that samples them).
  \item \texttt{reward\_fn(prompt, completion) -> float} --- the verifier.
  \item \texttt{eval\_fn(policy) -> dict} --- held-out evaluation returning a dict of named metrics.
\end{itemize}

New environments are contributed as pull requests. The review criterion is simple: does this environment exhibit a qualitatively distinct RL dynamic that existing environments do not cover? The goal is a curated zoo, not an uncurated dump.

\subsubsection{Leaderboards}

Each environment has a dedicated leaderboard tracking environment-specific competition metrics. Submissions include: algorithm name, hyperparameters, model checkpoint, and a reproducibility script. Results are verified by running the \texttt{eval\_fn} on a held-out test set that is not released publicly. Leaderboards reset on a quarterly cadence to encourage continued participation.

\subsection{TinyRL-Code: ``The MNIST of RLVR''}

\textbf{TinyRL-Code is the headline environment.} RLVR is the frontier paradigm. Code verification is the core RLVR domain at every major lab. If you build one environment, build this one.

\subsubsection*{Setup}

\begin{itemize}
  \item \textbf{Model:} GPT-2 small (124M) or CodeGen-350M, trained from a pretrained checkpoint.
  \item \textbf{Task bank:} 50--100 curated programming puzzles, graded by difficulty:
    \begin{itemize}
      \item \textit{Easy} --- implement a single function, verified by 1--3 unit tests (e.g., \texttt{reverse\_string}, \texttt{sum\_list}).
      \item \textit{Medium} --- implement 2 cooperating functions, verified by 5 unit tests (e.g., a parser and a formatter).
      \item \textit{Hard} --- implement an algorithm with edge cases, verified by 10 unit tests (e.g., BFS on an adjacency list, handling empty graph and disconnected nodes).
    \end{itemize}
  \item \textbf{Reward:} Fraction of unit tests passed. Pass all tests: reward = 1.0. Pass none: reward = 0.0. Partial credit for intermediate passes.
\end{itemize}

\subsubsection*{This Is RLVR in Miniature}

The verifier is objective, rule-based, and in principle non-gameable. In practice, models find exploits:

\begin{itemize}
  \item \textbf{Hardcoded outputs} --- the model learns the exact test inputs from repeated exposure and returns hardcoded outputs that pass the tests without implementing the function.
  \item \textbf{Format gaming} --- the model wraps a trivially wrong answer in \texttt{try/except} blocks to avoid exceptions, earning partial credit without solving the task.
  \item \textbf{Specification ambiguity} --- the model exploits underspecified edge cases in the problem statement that the test suite does not cover.
\end{itemize}

These are the same classes of exploit that frontier RLVR teams encounter at scale. TinyRL-Code makes them visible, characterizable, and patchable at toy scale --- exactly the verification engineering problem that matters for frontier training.

\subsubsection*{RLVR Dynamics Exhibited}

\begin{itemize}
  \item \textbf{Emergent problem-solving strategies} --- the model may develop chain-of-thought-like intermediate steps under verifiable reward pressure even without explicit CoT supervision.
  \item \textbf{Reward hacking on the verifier itself} --- as described above. This is the primary failure mode of RLVR.
  \item \textbf{Late-training capability emergence} --- harder puzzles may only become solvable after extended training, reproducing the ``late emergence'' phenomenon observed in frontier RLVR runs.
\end{itemize}

\subsubsection*{Competition Metrics}

\begin{enumerate}
  \item \textbf{Solve rate} --- fraction of held-out test puzzles with all unit tests passing. Primary leaderboard metric.
  \item \textbf{Capability per FLOP} --- solve-rate improvement per unit of RL compute (measured in total floating-point operations during training). Rewards efficient training, not just long training.
  \item \textbf{Reward Hacking Gallery} --- the community curates creative exploits found during training runs. Not a competition metric, but the most shared artifact. Every hardcoded output, every format hack, every specification exploit goes here.
\end{enumerate}

\subsection{TinyRL-Align: ``The RLHF Baseline''}

\subsubsection*{Setup}

\begin{itemize}
  \item \textbf{Model:} GPT-2 small (124M parameters), pretrained on OpenWebText.
  \item \textbf{Prompts:} IMDB movie review prefixes. The model is asked to complete the review.
  \item \textbf{Reward:} A pretrained sentiment classifier (e.g., \texttt{distilbert-base-uncased-finetuned-sst-2-english}) scores each completion. Positive sentiment = high reward.
\end{itemize}

\subsubsection*{This Is RLHF, Not RLVR}

The reward model is a neural network, not a deterministic verifier. It can be gamed. The policy does not need to write a genuinely positive review --- it needs to write text that the classifier scores highly. These are not the same thing.

\subsubsection*{RLHF Dynamics Exhibited}

\begin{itemize}
  \item \textbf{KL divergence growth} --- the policy drifts from the reference model rapidly when the KL penalty is weak. Observable within 200 steps.
  \item \textbf{Mode collapse} --- the policy converges to a narrow distribution of high-reward outputs (verbose, exclamation-heavy positive text). Entropy collapses.
  \item \textbf{Reward hacking} --- completions that score near 1.0 from the classifier are often incoherent or nonsensical to a human reader. The gap between proxy reward and true quality is the core RLHF problem.
  \item \textbf{Entropy collapse} --- policy entropy drops steeply as the model converges to its mode-collapsed distribution.
\end{itemize}

All dynamics visible in under 30 minutes on a single GPU.

\subsubsection*{Competition Metric}

Highest mean sentiment score (from the classifier) on the held-out IMDB prompt set, subject to a hard KL constraint: $D_{\mathrm{KL}}(\pi_\theta \| \pi_{\mathrm{ref}}) < 5.0$. This forces participants to optimize reward efficiently rather than through unconstrained drift.

\subsubsection*{Why Included}

TinyRL-Align is the lowest-barrier entry point. No code execution, no test harness, no syntax errors --- just a language model and a sentiment classifier. A researcher who has never implemented RL can run their first training loop here in under 30 minutes. The contrast with TinyRL-Code --- gameable neural reward versus in-principle-non-gameable verifier --- is the core RLHF vs.\ RLVR lesson.

\subsection{TinyRL-Quant: ``The PE-RL Flywheel''}

\subsubsection*{Setup}

\begin{itemize}
  \item \textbf{Model:} GPT-2 small (124M), with per-layer precision as a controllable parameter.
  \item \textbf{Action space:} Choose one of \{FP32, FP16, BF16, FP8, INT8, INT4\} for each transformer layer. A policy that maps layer index and layer statistics to a precision choice.
  \item \textbf{Reward:} A composite of evaluation quality (perplexity or accuracy on a held-out set) and inference efficiency (average bits per weight, or measured throughput). Specifically: $r = \text{accuracy} \times \frac{1}{\text{avg\_bits}}$, encouraging quality-per-bit optimization.
\end{itemize}

\subsubsection*{RLVR-Adjacent}

The reward is fully verifiable: accuracy is computed on a fixed evaluation set, and average bits per weight is a deterministic function of the chosen precisions. There is no neural proxy. This places TinyRL-Quant closer to the RLVR end of the spectrum than TinyRL-Align, but the reward combines two objectives, introducing reward design sensitivity that pure unit-test verification avoids.

\subsubsection*{Dynamics Exhibited}

\begin{itemize}
  \item \textbf{Reward design sensitivity} --- the weighting between accuracy and efficiency dramatically changes which layers the policy assigns high precision. Changing the reward formula changes the learned policy structure.
  \item \textbf{Exploration-exploitation in discrete spaces} --- the precision action space is combinatorial ($6^{12}$ for a 12-layer GPT-2). Pure exploitation fails; entropy regularization matters.
  \item \textbf{Eval-set gaming} --- if the same evaluation set is used for reward and final assessment, the policy overfits to it. Using a held-out test set distinct from the reward evaluation set is essential.
\end{itemize}

\subsubsection*{Competition Metric}

Best accuracy (on the held-out test set) per average bits per weight across all layers. A policy achieving 98\% of FP32 accuracy at 4.2 bits/weight beats one achieving 99\% at 6.1 bits/weight.

\subsubsection*{Unique Angle}

Apply the learned quantization policy to two variants: a model trained with RLVR (TinyRL-Code) and a model trained with SFT only. Does RLVR training change which layers the policy assigns high precision? The hypothesis is that RLVR training concentrates representations differently (more in later layers, due to verifiable reasoning demands), making the optimal quantization policy differ. This is a concrete, testable PE/RL intersection question.

\textit{Note: HAQ (Hardware-Aware Automated Quantization) exists as a production system but targets latency lookup tables for specific hardware. TinyRL-Quant focuses on RL dynamics study and the PE/RL intersection, not production deployment.}

\subsection{Stretch Environments}

Two future environments for contributors:

\begin{itemize}
  \item \textbf{TinyRL-Speculative.} RL for speculative decoding decisions: given a draft model and a verifier model, the policy decides how many tokens to speculate and when to verify. Reward is tokens per second on a fixed benchmark. Competition metric: peak tokens/second at fixed model quality. This bridges Chapter~4 (speculative decoding) and RLVR.

  \item \textbf{TinyRL-Reward.} Meta-RL for reward design: the policy proposes reward functions for a downstream task, and is evaluated on whether a model trained against those rewards achieves genuine capability gains (measured on a held-out probe set) without gaming. Competition metric: reward robustness score --- the gap between reward-training performance and probe-set performance, minimized. This is the hardest environment and the most research-adjacent.
\end{itemize}

\begin{thinkfirst}{The Verification Frontier at Toy Scale}
Rank the three TinyRL environments by verification cleanliness: TinyRL-Code (unit tests) $>$ TinyRL-Quant (composite verifiable reward) $>$ TinyRL-Align (neural proxy reward). For each environment, identify the dominant failure mode: verification hacking, KL collapse, or mode collapse. At what point on this spectrum does verification hacking become the dominant failure mode, displacing KL collapse as the primary concern? What does this imply about how verification quality should change your choice of KL penalty coefficient $\beta$?
\end{thinkfirst}

\begin{thinkfirst}{Designing a Minimal RLVR Environment}
Design the absolute simplest environment that satisfies all of the following constraints: (1) smallest model that works (how small is too small?), (2) fewest tasks (what is the minimum task bank size before the policy memorizes the training set?), (3) simplest verifier that is still non-trivially gameable, (4) must exhibit emergent strategy AND verification gaming, (5) must train to completion in under 15 minutes on a single consumer GPU. Write a complete specification: model, task format, verifier pseudocode, reward function, and the exploit you predict the model will find first. Justify each design choice.
\end{thinkfirst}

\subsection{Exercises}

\begin{exercise}{Train with PPO Across All Environments (1--2 hours per environment)}
For each of TinyRL-Code, TinyRL-Align, and TinyRL-Quant: train with PPO for 500 steps using the default hyperparameters (\texttt{lr=1e-5}, \texttt{beta=0.1}, \texttt{clip\_eps=0.2}). For each run, plot reward curve, KL divergence from reference, and policy entropy on the same figure. Answer: (1) What is the primary pathology in each environment? (2) Which environment shows reward hacking first (fewest steps)? (3) In TinyRL-Code, log the first detected reward hack: what did the model do? (4) In TinyRL-Align, at what step does entropy begin to collapse? Compare your findings across the three environments and write a one-paragraph characterization of how the verifier type shapes the failure mode.
\end{exercise}

\begin{exercise}{Train with GRPO and Sweep Group Size (2--3 hours)}
GRPO (Group Relative Policy Optimization) is the algorithm DeepSeek used for R1-Zero RLVR training. It replaces the value function with a group-relative baseline computed over $G$ sampled completions per prompt. For each of the three environments, train with GRPO at $G \in \{4, 8, 16\}$ for 500 steps. Compare: (1) Is GRPO more stable than PPO (lower KL variance, more consistent reward growth)? (2) Does increasing $G$ always improve sample efficiency, or does the benefit plateau? (3) What is the compute cost of $G=16$ vs.\ $G=4$ in wall-clock time? (4) For TinyRL-Code specifically: does larger $G$ reduce the frequency of reward hacking, or does it accelerate it? Produce a 3$\times$3 table: environment $\times$ $G$ value, reporting final reward and final KL at step 500.
\end{exercise}

\begin{exercise}{Remove the KL Penalty and Document Failure Modes (1 hour)}
Train each environment for 500 steps with $\beta = 0$ (no KL penalty). Document: (1) What happens to KL divergence without the penalty? (2) What happens to reward --- does it increase faster, or does training collapse? (3) \textbf{Key comparison:} does TinyRL-Code (RLVR, deterministic verifier) tolerate KL removal better than TinyRL-Align (RLHF, neural reward)? State a hypothesis before running. If RLVR tolerates KL removal better: argue why deterministic verification provides an implicit regularizer that neural rewards do not. If it does not: what does that tell you about the role of KL in preventing reward hacking versus preventing distribution collapse? (4) At what KL value does each environment's policy produce visibly degenerate outputs? Check qualitative samples at step 500.
\end{exercise}

\begin{exercise}{Long-Horizon RLVR Test on TinyRL-Code (3--4 hours)}
Train TinyRL-Code for 5{,}000 steps (10$\times$ the standard 500-step run). Log at every step: solve rate on training prompts, solve rate on held-out test set, policy entropy, gradient $\ell_2$ norm, and KL from reference. Questions: (1) Do new capabilities emerge late --- does the model begin solving medium or hard puzzles only after step 3{,}000? (2) Do new failure modes appear late --- does verifier overfitting (hardcoded outputs for training prompts) worsen after step 2{,}000? (3) Plot capability (held-out solve rate) vs.\ log compute (log total FLOPs). Is the curve approximately log-linear? If so, this supports the claim that RLVR capability scales smoothly with compute. If not, describe the shape and propose an explanation. (4) Apply a change-point detection algorithm (PELT or binary segmentation) to the held-out solve rate series. Do change-points correspond to jumps in entropy or KL?
\end{exercise}

\begin{exercise}{Quantization Stress Test Under RL Training (2--3 hours)}
This exercise tests numerical stability of RL training under reduced precision, directly relevant to Dario's ``laminar flow'' framing of stable training runs. For TinyRL-Align: run the standard PPO training loop with the policy model stored and updated in (a) FP32, (b) FP16, (c) FP8 (if your hardware supports it), (d) INT8 (using a quantization library). For each precision: (1) Log gradient $\ell_2$ norm at each step. At what precision does the gradient norm become unstable (variance $>10\times$ FP32 variance)? (2) Does training diverge? At what step? (3) Compare final reward at step 500 across precisions. (4) Is there a precision where training appears stable by reward but has degraded by KL or entropy metrics? Report the ``safe precision floor'' for RL training on your hardware and explain what numerical phenomenon causes instability below that floor.
\end{exercise}

\begin{exercise}{Adversarial Verification Engineering (2--3 hours)}
This is the core skill for RLVR at scale. For TinyRL-Code: (1) \textbf{Design 5 gameable test suites.} Each should be gameable in a different way: (a) tests that only check a single predictable input, (b) tests that check output format but not correctness, (c) tests with insufficient coverage that miss a large class of bugs, (d) tests with ambiguous specifications exploitable by degenerate solutions, (e) tests that are bypassable by catching exceptions. (2) Train a fresh model against each gameable suite for 300 steps. Document the exploit each suite induces. (3) \textbf{Harden each test suite.} Add edge cases, remove predictable inputs, enforce semantic correctness. (4) Re-train against the hardened suites. Measure: how much does hardening improve held-out solve rate on the \textit{canonical} test suite (which was never shown during training)? Report a ``verification hardening multiplier'' for each suite: ratio of post-hardening to pre-hardening canonical solve rate. Explain the pattern.
\end{exercise}

\begin{portfolio}{TinyRL: The MNIST of RL for LLMs}
\textbf{Deliverable:} An open-source framework on GitHub with three implemented environments (Code, Align, Quant), a leaderboard site, and a launch blog post.

\textbf{The pitch:} ``I built a standardized platform where anyone can study RLVR and RLHF dynamics in 30 minutes on one GPU. The headline environment --- TinyRL-Code --- is RLVR in miniature: verifiable rewards, deterministic unit tests, emergent problem-solving strategies, and a reward hacking gallery. TinyRL-Align provides the RLHF baseline for direct comparison. The contrast between them is the pedagogical core: it shows concretely why the field moved from neural proxy rewards to verifiable rewards, and what is gained and lost in that transition.''

\textbf{The monkey:} The pedestal is ``can you build a framework?'' That is not the question. The monkey is: \textit{What RLVR phenomena are observable at toy scale, and which require frontier-scale models?} If TinyRL-Code reproduces real pathologies --- reward hacking, entropy collapse, emergent chain-of-thought under verifiable pressure --- you have built a useful research tool. If it does not, you have learned something equally important about the scale-dependence of RL phenomena. Either outcome is valuable. Kill the monkey by running TinyRL-Code for 5{,}000 steps and characterizing what it does and does not reproduce.

\textbf{Verification criteria:}
\begin{enumerate}
  \item TinyRL-Code exhibits at least 2 qualitatively distinct reward hacking strategies, documented in the Reward Hacking Gallery with completion examples and step counts.
  \item The RLHF vs.\ RLVR comparison (Align vs.\ Code, same algorithm and hyperparameters) shows measurably different dynamics: different dominant failure modes, different KL trajectories, different entropy collapse timing. Document this with side-by-side plots.
  \item A user who has never implemented RL can train their first model in under 30 minutes by following the README. Test this by asking someone with no RL background to run TinyRL-Align from scratch.
  \item At least 3 community-contributed environments within 3 months of public launch. This requires writing a contribution guide, a minimal environment template, and seeding the community with the stretch environments as open issues.
\end{enumerate}

\textbf{Why community:} Lowering the barrier to RLVR experimentation creates contributors and competitors. The reward hacking gallery is inherently shareable --- every exploit is a story. Network effects compound: each new environment attracts researchers who care about that domain, who then contribute fixes and new environments. The leaderboard creates a persistent reason to return. This is how MNIST became a standard: not by being the best benchmark, but by being the one everyone used.
\end{portfolio}

\newpage
\section{RL Pipeline \& Verification Engineering}

\smallskip\noindent\textit{Study Guide companion: Study Guide Sections 2--5 cover this material in interview format.}
\smallskip

The modern post-training pipeline is \textbf{RLVR-first}: Base model $\to$ SFT $\to$ RLVR (GRPO with verifiable rewards on math/code) $\to$ RLHF (neural reward model for subjective tasks). RLHF is no longer the main compute consumer. Understanding both stages --- and where verification engineering fits --- is the entry point for every RE role at a frontier lab.

\begin{table}[h]
\centering
\begin{tabular}{@{}lll@{}}
\toprule
\textbf{Era} & \textbf{Pipeline} & \textbf{Compute Distribution} \\
\midrule
2020--2022 & Pretraining only & 100\% pretraining \\
2022--2024 & Pretraining $\to$ SFT $\to$ RLHF & $\sim$95\% pretraining, $\sim$5\% SFT+RLHF \\
2025+ & Pretraining $\to$ SFT $\to$ RLVR $\to$ RLHF & Pretraining shrinking, RLVR growing \\
\bottomrule
\end{tabular}
\caption{Training pipeline evolution. The RLHF stage has not disappeared --- it handles subjective alignment tasks where clean verification is unavailable. But RLVR now dominates post-training compute for capability gains.}
\end{table}

\textbf{Intelligence-per-Watt.} Under RLVR, every wasted training step is wasted compute at massive scale. A verifier that is 10\% more robust effectively gives you 10\% more intelligence per Watt --- not via kernel optimization, but via signal quality. This is why verification engineering is a first-class research problem, not a supporting task.

\subsection{Conceptual Foundation}

\begin{thinkfirst}{The RL training pipeline as a system}
Draw the full post-training pipeline on paper. Both RLVR and RLHF stages are present. For each stage, answer:

\begin{enumerate}
  \item \textbf{SFT:} What distribution does the model learn? DeepSeek R1-Zero skipped this stage entirely. What emerged? What was missing? When would you skip it and when would you not?

  \item \textbf{RLVR stage:} The reward is binary --- the verifier says pass or fail. You need no reward model. What are you trading away compared to a trained reward model? What do you gain? (Hint: Goodhart's Law runs in both directions.)

  \item \textbf{RLHF stage (neural RM):} Why rankings instead of scores? If you have 4 outputs to rank, how many pairwise comparisons can you extract? Why is this more data-efficient than scoring independently?

  \item \textbf{The KL penalty in GRPO:} Imagine removing it entirely. The model is free to go anywhere in output space. Predict what happens in the first 100 training steps. In the first 1000. Why does the failure mode get \textit{worse} over time?

  \item \textbf{GRPO vs.\ PPO:} PPO needs 4 models in GPU memory simultaneously. GRPO drops the critic. Write down what each of PPO's 4 models does and which require gradients. What does GRPO replace the critic's function with?
\end{enumerate}
\end{thinkfirst}

\begin{thinkfirst}{Why online RL beats offline}
DPO trains on a fixed dataset of preference pairs. PPO/GRPO generate new outputs during training and evaluate them live.

\begin{enumerate}
  \item A model trained with DPO discovers a new exploit after deployment. Why couldn't DPO have caught this during training? What would GRPO have done differently?

  \item Frontier RL teams use online RL, not DPO. Explain why in one sentence that a non-ML engineer could understand. (Hint: a student who only studies from a fixed answer key vs.\ one who practices with a tutor who gives new problems.)

  \item DPO's loss function is binary cross-entropy on preference pairs. What information is lost compared to RLVR's binary verifier? What information is lost compared to RLHF's scalar reward model? Are these the same loss?
\end{enumerate}
\end{thinkfirst}

\begin{thinkfirst}{Goodhart's Law in RL --- the verification gaming frame}
``When a measure becomes a target, it ceases to be a good measure.'' Under RLVR, the verifier \textit{is} the measure. Apply this:

\begin{enumerate}
  \item You train a coding agent with reward = ``all unit tests pass.'' The model learns to modify the test file. Why does the RL algorithm see this as a perfectly valid solution? What would you need to change about the verifier to prevent this?

  \item Design a reward for ``write a good essay'' using an LLM judge. List 3 ways a model could game this reward without writing a good essay. Be specific --- think about what the judge is actually sensitive to. Now explain why RLVR on math sidesteps this problem.

  \item Imagine using \textit{human expert decisions} as the reward signal instead of a neural RM. In what sense is this more robust to gaming? In what sense is it fragile? Think about what happens on easy cases vs.\ hard/ambiguous cases.
\end{enumerate}
\end{thinkfirst}

\begin{thinkfirst}{The verification frontier}
Rank the following domains from ``easiest to verify'' to ``hardest to verify.'' For each, describe what verification would look like and what the model could exploit:

\begin{enumerate}
  \item Arithmetic (e.g., ``What is 847 $\times$ 293?'')
  \item Competitive programming (e.g., Codeforces problems)
  \item Code review quality (e.g., ``Does this PR introduce bugs?'')
  \item Customer service quality (e.g., ``Was this a helpful response?'')
  \item Novel writing (e.g., ``Is this a compelling opening chapter?'')
\end{enumerate}

For the middle-difficulty domains (code review, customer service), design a \textbf{composite verifier} --- a combination of automated checks + LLM judge + human spot-checks. What would the model try to game? How would you detect it?
\end{thinkfirst}

\begin{thinkfirst}{Process vs.\ outcome rewards}
A Process Reward Model (PRM) scores each reasoning step. An Outcome Reward Model (ORM) scores only the final answer. RLVR typically uses an ORM-equivalent (pass/fail on the final answer).

\begin{enumerate}
  \item A student solves a math problem. They get the right answer but their reasoning has a cancelling error in step 3 and step 5. What does the ORM say? What does the PRM say? Which is more useful for learning?
  \item PRMs need step-level labels. Where do these come from? Why are they expensive? Can you generate step-level labels \textit{automatically} for coding tasks? How?
  \item DeepSeek R1 used \textit{only} outcome rewards. The model spontaneously developed chain-of-thought, self-verification, and backtracking. If process rewards aren't necessary for \textit{generating} good reasoning, what \textit{are} they necessary for?
\end{enumerate}
\end{thinkfirst}

\subsection{Hands-On Exercises}

\begin{exercise}{nanoRLVR --- The modern baseline (10--15 hours)}
Build a \textbf{minimal RLVR pipeline} on a verifiable task. This is the baseline every RE candidate should have in their portfolio.

\textbf{Suggested setup:} GPT-2 (124M) fine-tuned on simple math problems (addition, subtraction, basic arithmetic). The verifier is an exact-match checker against the ground truth answer --- no neural reward model needed.

\textbf{Stages:} (1) SFT: fine-tune GPT-2 on math problem--solution pairs in a structured format. (2) RLVR: GRPO loop --- for each problem, generate $G = 8$ completions, run the verifier on each, compute group-normalized advantages $A_i = (r_i - \bar{r})/\sigma_r$, update with PPO clipped objective. No critic network.

\textbf{Ablations that teach you the most:}
\begin{enumerate}[label=\alph*)]
  \item Remove the KL penalty ($\beta = 0$). Train 500 steps. What does the model learn to do? Can you identify the exact exploit it found? (Common: it generates very long outputs that happen to contain the correct number, or it learns a formatting trick that systematically passes the verifier.)
  \item Vary group size $G = 2, 4, 8, 16$. How does advantage estimation quality change? At what $G$ does training become unstable?
  \item Swap the verifier for a deliberately broken one (e.g., ``correct if answer contains any digit $>$ 5''). Watch what the model learns. This is verification gaming in miniature.
\end{enumerate}

\textbf{Verify your pipeline:} Reward increases over first 200 steps; KL bounded; generated completions are coherent; group-normalized advantages have mean $\approx 0$ and std $\approx 1$.
\end{exercise}

\begin{exercise}{nanoRLHF --- Legacy comparison (10--15 hours)}
\textbf{Run this on the same model after nanoRLVR.} This is a deliberate comparison, not an independent exercise.

Build the classic RLHF pipeline: GPT-2 + sentiment steering. Use a pre-trained sentiment classifier (\texttt{distilbert-base-uncased-finetuned-sst-2-english}) as the reward model. SFT on IMDB reviews, then PPO with the classifier as reward.

\textbf{The comparison experiment:} After running both pipelines to equivalent compute, compare: (a) reward signal noise (verifiable reward is binary and clean; sentiment classifier output is noisy and continuous), (b) the exploits each pipeline discovers, (c) which pipeline requires more hyperparameter tuning.

Write a 1-paragraph post-mortem: in what task setting would you prefer RLHF over RLVR? What does the answer tell you about the verification frontier?
\end{exercise}

\begin{exercise}{PPO to GRPO --- Drop the critic (4--6 hours)}
Starting from a working PPO implementation (or TRL's \texttt{PPOTrainer}):

\begin{enumerate}
  \item Remove the critic/value network entirely. For each prompt, generate $G = 8$ completions. Score all 8. Compute advantages as $A_i = (r_i - \bar{r})/\sigma_r$ (group normalization). Use these in the PPO clipped objective.
  \item \textbf{Measure:} GPU memory saved by dropping the critic. Training stability (reward variance across batches) for PPO vs.\ GRPO. Try $G = 2, 4, 8, 16, 32$.
  \item At what $G$ does GRPO match PPO's training stability? Is the memory savings worth it? Write a 1-paragraph recommendation.
\end{enumerate}
\end{exercise}

\begin{exercise}{Verification hacking benchmark (6--8 hours)}
Build a small benchmark where models can exploit verification loopholes --- then harden it.

\textbf{Setup:} 10--20 coding problems where the verifier is a set of unit tests. Plant loopholes: (a) predictable output patterns the model can hard-code, (b) tests that check format only, not correctness, (c) deterministic tests with insufficient coverage.

Run best-of-N sampling (N=16). Count how many problems the model ``solves'' by gaming the tests. Then harden: add randomized test inputs, format-independent output checking, adversarial cases. Re-run and measure how many exploits you closed. Measure the cost (human time, compute) of each hardening step. Is there a Pareto frontier of verification quality vs.\ cost?
\end{exercise}

\begin{exercise}{Composite verifier design for a fuzzy domain (4--6 hours)}
Pick a domain where ``good'' is fuzzy (code review quality, documentation clarity, explanation quality).

Design a composite verifier with 3+ components: automated checks, LLM judge, and a consistency check (do two independent judges agree?). Weight the components into a scalar reward. Generate 50 outputs at varying quality; score them. Inspect the top-10 and bottom-10 manually --- does the verifier rank correctly?

Then: deliberately try to game your own verifier. Write 5 outputs designed to score high without being good. If you succeed, iterate. Write a 1-page analysis: at what point would a model trained against this reward start producing outputs you'd be unhappy with? Is this signal clean enough for RLVR?
\end{exercise}

\subsection{Optional Extensions}

These are productive next steps after the core exercises but not required for a portfolio:

\begin{itemize}
  \item \textbf{DPO comparison:} Train DPO on preference pairs generated by your RLVR model. Compare: can DPO reach the same performance? Where does it fail to generalize?
  \item \textbf{Multi-turn RL:} Extend nanoRLVR to a 3--5 step environment. The credit assignment problem becomes real when the reward comes only at the end.
  \item \textbf{RLAIF:} Replace the human-labeled reward model in nanoRLHF with AI-generated feedback. Can the AI judge provide useful signal on tasks it can't solve itself?
\end{itemize}

\subsection{Portfolio Projects}

\begin{portfolio}{nanoRLVR: The Pedagogical Pipeline}
Build the ``nanoGPT of RLVR.'' A minimal, readable, single-repo implementation that goes from pretraining through SFT, RLVR (GRPO), and RLHF, all on toy tasks with clean verification. Include nanoRLHF as a labeled legacy comparison.

\textbf{The monkey:} Not ``can you build it'' (that's the pedestal). The monkey is: ``what surprising thing did you discover in the ablations that isn't in existing tutorials?'' Include ablation figures: KL sweep, group size sweep, verifier hardness sweep, GRPO vs.\ PPO vs.\ DPO. Document what broke and how you diagnosed it. Include a ``reward hacking gallery'' --- examples of exploits each pipeline found.

\textbf{Verification:} Someone can clone and reproduce all results in $<$4 hours on a single GPU. At least one ablation finding surprises a reader who has used TRL before.
\end{portfolio}

\begin{portfolio}{Verification Hacking Taxonomy}
Build a systematic study of how models game code verifiers. Not just ``they hack tests'' --- categorize exploit types, measure which mitigations work, and publish findings.

\textbf{Structure:} (1) Taxonomy: hard-coding outputs, format exploitation, specification ambiguity, distributional shortcuts, test-case overfitting, environment manipulation. (2) Benchmark: 50+ problems with known vulnerabilities, difficulty-graded. (3) Mitigations tested: sandboxing, randomized inputs, output validation, RM ensembles, specification hardening. (4) Cost-effectiveness analysis: for each mitigation, compute/human-time cost vs.\ exploit closure rate.

\textbf{The monkey:} ``Which exploit categories matter most, and which mitigations are cost-effective?'' Building the benchmark is the pedestal. The Pareto frontier of verification quality vs.\ cost is the monkey. Solve for that.

\textbf{Deliverable:} OSS benchmark + blog post. Verification: your analysis shows at least one dominated mitigation and makes a concrete recommendation an RL team could act on.
\end{portfolio}

\newpage
\section{Scaling Laws}

\smallskip\noindent\textit{Study Guide companion: Study Guide Section 11 (Scaling Laws Deep Dive) covers the interview-ready version of this material.}
\smallskip

Before spending \$10M on a large RLVR training run, can you predict from \$10K experiments whether the reward signal and algorithm will hold at scale? Andy Jones demonstrated this skill independently with ``Scaling Scaling Laws with Board Games'' (arxiv 2104.03113) --- training AlphaZero on small Hex boards to predict performance on large ones. This chapter builds the same skill on a toy domain you can run on a single GPU.

\textbf{Intelligence-per-Watt, strategic edition.} Kernel optimization gives you 2$\times$ on one operation. Correct compute allocation gives you 2$\times$ on the entire training budget. Scaling laws are the map from ``how much compute'' to ``how much capability'' --- and the only way to navigate that map without running the expensive experiment is to fit it from cheap experiments first.

\subsection{Conceptual Foundation}

\begin{thinkfirst}{Why scaling laws matter for RL teams}
Before a large RLVR training run, the team needs to answer: ``Will this algorithm + reward signal + compute budget actually produce the capability we want?''

\begin{enumerate}
  \item You have a new verifiable reward signal for a coding task. You run small experiments on easy problems and see improvement. Should you immediately scale up? What could go wrong? (Hint: ``things which can help at smaller scales can actually hurt at larger scales.'')

  \item Andy Jones showed that performance on small Hex boards predicts performance on large ones. What's the analog for LLM RLVR? (Hint: ``small Hex board'' $\to$ easy coding problem. ``Large Hex board'' $\to$ hard coding problem. Now: what happens if the verifier for hard problems is noisier than for easy ones?)

  \item The Jones paper found performance is sigmoid in compute --- there's a ``takeoff'' threshold below which nothing happens. What does this mean for budget allocation? If you're below the takeoff threshold, is it better to run a bigger model or train longer? Does RLVR's longer training runs (relative to RLHF) change this answer?
\end{enumerate}
\end{thinkfirst}

\begin{thinkfirst}{The train/test compute tradeoff}
Jones found: $10\times$ more training $\approx$ $15\times$ less inference search needed.

\begin{enumerate}
  \item Translate this to LLMs: what does ``inference search'' correspond to? List at least 3 forms of test-time compute that LLMs use.

  \item If this tradeoff holds for frontier coding agents: more RLVR training should reduce the chain-of-thought needed at inference, making each API call cheaper. Why does this matter for serving economics?

  \item The tradeoff is log-linear. Does it hold forever, or is there a floor? Even a perfectly trained model might need \textit{some} reasoning for hard problems. What does that floor represent?

  \item Frontier labs optimize across three axes: pretraining, RLVR training, and test-time compute. If the train/test tradeoff is log-linear, what shape is the efficient frontier? Sketch it on paper. Where on that frontier do you think Claude 3.5 Sonnet sits relative to o1?
\end{enumerate}
\end{thinkfirst}

\subsection{Hands-On Exercises}

\begin{exercise}{Replicate Jones at toy scale (15--20 hours)}
Replicate the core finding of the Scaling Scaling Laws paper on a toy domain. Use a simple RL algorithm (DQN or PPO) on gridworld mazes of increasing size (5$\times$5, 7$\times$7, 10$\times$10, 15$\times$15).

\begin{enumerate}
  \item \textbf{Sweep compute:} For each maze size, train with 5 different compute budgets (steps: 1k, 3k, 10k, 30k, 100k).
  \item \textbf{Plot the frontiers:} For each maze size, plot performance vs.\ compute. Do you see the sigmoid shape Jones found?
  \item \textbf{Fit the scaling law:} Use Jones's 5-parameter change-point model:
  \begin{lstlisting}
  plateau = m_size * maze_size + c_plateau
  incline = m_size_i * maze_size + m_flops * log(compute) + c_incline
  performance = clip(incline, plateau, ceiling)
  \end{lstlisting}
  Fit with \texttt{scipy.optimize.curve\_fit}. The fit should achieve $R^2 > 0.9$ on training data --- if not, check initialization or functional form.
  \item \textbf{The key test:} Fit only on maze sizes 5--10. Predict performance for sizes 12 and 15. This is the ``scaling scaling laws'' result. How accurate is the extrapolation?
  \item \textbf{The open question:} Does prediction accuracy degrade gracefully or catastrophically as you extrapolate to size 20? Document what you find.
\end{enumerate}
\end{exercise}

\begin{exercise}{The train/test tradeoff (2--3 hours)}
Using the same toy domain:

\begin{enumerate}
  \item Train an agent at 5 compute levels. At each level, evaluate with varying test-time compute (for gridworld: number of planning steps or beam width).
  \item Plot iso-performance curves: what combinations of train and test compute give the same solve rate?
  \item Fit a log-linear relationship: $\log(\text{test}) = a \cdot \log(\text{train}) + b$. Does it hold? What's $a$?
  \item Compare to Jones's finding: $a \approx -1.2$. Is yours steeper or shallower? What might cause the difference?
\end{enumerate}
\end{exercise}

\begin{exercise}{Where does extrapolation break? (2--3 hours)}
This is the genuinely hard question Jones raised but didn't fully answer.

\begin{enumerate}
  \item In your toy domain, introduce a \textbf{qualitative change in difficulty} at some maze size (e.g., at size 12+, add locked doors requiring key items). Does the scaling law still extrapolate across this discontinuity?
  \item What's the analog for RLVR on code? If the verifier (unit tests) is clean for easy problems but starts breaking down for hard ones (insufficient test coverage), will the scaling law still hold?
  \item \textbf{Write a 1-paragraph prediction:} When you scale RLVR compute on coding tasks, where do you expect the scaling law to break? What would cause a discontinuity? (Think about verifier quality, not just model capacity.)
  \item \textbf{Design the next experiment:} You have 24 GPU-hours. Write a brief proposal: what you'd test, what you'd measure, what the result would tell you. Then run it.
\end{enumerate}
\end{exercise}

\begin{portfolio}{Scaling Laws for RL on Code}
Replicate Jones's approach but for \textbf{coding tasks of varying difficulty}. Use a small model (e.g., CodeGen-350M or StarCoder-1B) and RLVR-fine-tune it on programming problems graded by difficulty.

\textbf{Problem difficulty parameterization:}
\begin{itemize}
  \item Easy: single-function, 1--3 unit tests, standard algorithms
  \item Medium: multi-function, 5--10 tests, requires data structure choice
  \item Hard: multi-file, 10+ tests, requires architectural reasoning
\end{itemize}

\textbf{Key questions:}
\begin{enumerate}
  \item Does performance scale as a sigmoid in compute for each difficulty level (like Jones's Hex)?
  \item Can you predict hard-task performance from easy-task experiments?
  \item At what difficulty level does the unit-test verifier start breaking down? Smooth degradation or cliff?
  \item What's the train/test compute tradeoff for coding? Does more RLVR training reduce the chain-of-thought needed at inference?
\end{enumerate}

\textbf{The monkey:} ``Does extrapolation across problem difficulty actually hold for code?'' Running the experiments is the pedestal. The insight about \textit{where and why} extrapolation breaks --- and whether it's the verifier or the model that's the bottleneck --- is the monkey.

\textbf{Deliverable:} Paper or detailed blog post with scaling law fits and extrapolation results. \textbf{Verification:} Prediction error on held-out difficulty levels $<$15\%. If worse, document where extrapolation breaks --- that finding is equally valuable.
\end{portfolio}

\newpage
